\documentclass[11pt]{article}
\usepackage[utf8]{inputenc}
\usepackage{amsmath, amssymb}
\usepackage{geometry}
\geometry{margin=1in}

\title{Two-Sector DMP Model with Continuous Quality Sorting}
\author{Based on Mortensen-Pissarides (1994) and Lagos (2006)}
\date{\today}

\begin{document}

\maketitle

\section{Environment}
\begin{itemize}
    \item \textbf{Time:} Continuous.
    \item \textbf{Sectors:} $j \in \{H, L\}$ denoting High-tech and Low-tech (non-high-tech). Sector-specific productivity factors are $p_H$ and $p_L$.
    \item \textbf{Workers:} Unit mass of risk-neutral workers. Workers are heterogeneous in ability $a$, which is distributed continuously according to a CDF $F(a)$ with density $f(a)$ on support $[a_{\min}, a_{\max}]$.
    \item \textbf{Labor Force:} The total mass of workers with ability $a$ is $f(a)$. Workers discount the future at rate $r$. Unemployed workers with ability $a$ receive a flow benefit proportional to their ability, $b \cdot a$.
    \item \textbf{Firms:} Large pool of risk-neutral potential firms. A firm can open a vacancy in sector $j$ at a flow cost $c_j$.
\end{itemize}

\section{Search and Matching}
Workers direct their search to a specific sector based on expected utility, but match randomly with firms within that sector.
\begin{itemize}
    \item \textbf{Sector Choice:} Let $\phi_{j}(a) \in [0, 1]$ be the probability/fraction of unemployed workers of ability $a$ who choose to search in sector $j \in \{H, L\}$.
    \item \textbf{Sector Pools:} The mass of unemployed workers searching in sector $j$ is $u_j = \int \phi_{j}(a) u(a) da$, where $u(a)$ is the density of unemployed workers of ability $a$.
    \item \textbf{Market Tightness:} Let $v_j$ be the number of vacancies in sector $j$. Market tightness is $\theta_j = v_j / u_j$.
    \item \textbf{Contact Rates:} The matching function is $M(u_j, v_j)$. The worker job-finding rate is $\lambda(\theta_j) = M(u_j, v_j) / u_j$. The firm vacancy-filling rate is $q(\theta_j) = M(u_j, v_j) / v_j$.
    \item \textbf{Pool Composition:} Conditional on a firm in sector $j$ meeting a worker, the density of the worker's ability is $\pi_{j}(a) = (\phi_{j}(a) u(a)) / u_j$.
\end{itemize}

\section{Production and Endogenous Separation}
\begin{itemize}
    \item \textbf{Match Quality:} Upon meeting, the firm and worker draw a match-specific idiosyncratic quality shock $x \sim U[0, 1]$. The uniform distribution has mean $0.5$, density $g(x) = 1$, and CDF $G(x) = x$.
    \item \textbf{Output:} A producing match between a sector $j$ firm and a worker of ability $a$ with shock $x$ yields flow output $y_j(a, x) = p_j a x$.
    \item \textbf{Shocks:} New idiosyncratic shocks $x'$ arrive at Poisson rate $\delta$, drawn from the same uniform distribution $U[0, 1]$.
    \item \textbf{Exogenous Separation:} Matches are destroyed for exogenous reasons at rate $s$.
    \item \textbf{Endogenous Separation:} A match is continued if and only if the joint surplus is non-negative, defining a reservation match quality $R_{j}(a)$. If $x < R_{j}(a)$, the match separates.
\end{itemize}

\section{Value Functions}

\subsection{Workers}
The value of unemployment for a worker of ability $a$ searching in sector $j$ is $U_j(a)$:
\begin{equation}
r U_j(a) = b a + \lambda(\theta_j) \int_{R_j(a)}^1 [W_j(a, x) - U_j(a)] dx
\end{equation}

The value of employment $W_j(a, x)$ satisfies:
\begin{equation}
r W_j(a, x) = w_j(a, x) + \delta \int_{R_j(a)}^1 [W_j(a, x') - W_j(a, x)] dx' - (\delta R_j(a) + s) [W_j(a, x) - U_j(a)]
\end{equation}

\textbf{Optimal Sorting:} Workers choose the sector that maximizes their value:
\begin{equation}
U(a) = \max \{ U_H(a), U_L(a) \}
\end{equation}

\subsection{Firms}
The value of a vacancy in sector $j$ is $V_j$:
\begin{equation}
r V_j = -c_j + q(\theta_j) \int \pi_j(a) \int_{R_j(a)}^1 [J_j(a, x) - V_j] dx da
\end{equation}
Free entry implies $V_j = 0$ for all active sectors.

The value of a filled job $J_j(a, x)$ satisfies:
\begin{equation}
r J_j(a, x) = p_j a x - w_j(a, x) + \delta \int_{R_j(a)}^1 [J_j(a, x') - J_j(a, x)] dx' - (\delta R_j(a) + s) J_j(a, x)
\end{equation}

\section{Surplus and Equilibrium Conditions}
Wages are determined by continuous Nash Bargaining with worker bargaining power $\beta$:
\begin{equation}
W_j(a, x) - U_j(a) = \beta S_j(a, x) \quad \text{and} \quad J_j(a, x) - V_j = (1-\beta) S_j(a, x)
\end{equation}

The total match surplus $S_j(a, x) = W_j(a, x) - U_j(a) + J_j(a, x) - V_j$ satisfies:
\begin{equation} \label{eq:surplus}
(r + s + \delta) S_j(a, x) = p_j a x - r U_j(a) + \delta \int_{R_j(a)}^1 S_j(a, x') dx'
\end{equation}

\subsection{Job Destruction (Reservation Quality)}
The reservation quality $R_j(a)$ is defined by $S_j(a, R_j(a)) = 0$:
\begin{equation} \label{eq:reservation}
p_j a R_j(a) - r U_j(a) + \delta \int_{R_j(a)}^1 S_j(a, x') dx' = 0
\end{equation}

Subtracting \eqref{eq:reservation} from \eqref{eq:surplus}, we can express the surplus linearly in $x$:
\begin{equation} \label{eq:surplus_linear}
S_j(a, x) = \frac{p_j a (x - R_j(a))}{r + s + \delta}
\end{equation}

Substituting \eqref{eq:surplus_linear} into \eqref{eq:reservation} and evaluating the integral over the uniform distribution $x \sim U[0,1]$ yields:
\begin{equation} \label{eq:reservation_explicit}
p_j a R_j(a) - r U_j(a) + \frac{\delta p_j a}{r + s + \delta} \frac{(1 - R_j(a))^2}{2} = 0
\end{equation}

\subsection{Derivatives of the Surplus Function}
Using \eqref{eq:surplus_linear}, we can derive the partial derivatives of the surplus $S_j(a, x)$ with respect to match quality $x$ and ability $a$.

Differentiating with respect to $x$:
\begin{equation}
\frac{\partial S_j(a, x)}{\partial x} = \frac{p_j a}{r + s + \delta}
\end{equation}

Differentiating with respect to $a$:
\begin{equation} \label{eq:ds_da_full}
\frac{\partial S_j(a, x)}{\partial a} = \frac{p_j (x - R_j(a)) - p_j a R'_j(a)}{r + s + \delta}
\end{equation}

\textit{Note on proportionality:} Using the Nash bargaining solution, the value of unemployment can be written as:
\begin{equation} \label{eq:unemp_value}
r U_j(a) = b a + \lambda(\theta_j) \beta \int_{R_j(a)}^1 \frac{p_j a (x - R_j(a))}{r + s + \delta} dx = b a + \lambda(\theta_j) \beta \frac{p_j a (1 - R_j(a))^2}{2(r + s + \delta)}
\end{equation}
Equating \eqref{eq:unemp_value} with the expression for $r U_j(a)$ from \eqref{eq:reservation_explicit} shows that $a$ cancels out entirely from the equation determining $R_j(a)$. Therefore, $R_j(a) = R_j$ is independent of ability $a$, which implies $R'_j(a) = 0$.

Consequently, the partial derivative with respect to $a$ from \eqref{eq:ds_da_full} simplifies to:
\begin{equation}
\frac{\partial S_j(a, x)}{\partial a} = \frac{p_j (x - R_j)}{r + s + \delta}
\end{equation}

\subsection{Job Creation (Free Entry)}
Imposing $V_j = 0$ yields the job creation condition for sector $j$:
\begin{equation}
\frac{c_j}{q(\theta_j)} = (1-\beta) \int \pi_j(a) \int_{R_j(a)}^1 S_j(a, x) dx da = (1-\beta) \int \pi_j(a) \frac{p_j a (1 - R_j(a))^2}{2(r + s + \delta)} da
\end{equation}

\subsection{Steady-State Flow Equations}
Let $e_j(a)$ be the density of employed workers of ability $a$ in sector $j$.
\begin{equation}
u(a) + e_H(a) + e_L(a) = f(a)
\end{equation}
The flow into employment in sector $j$ must equal the flow out:
\begin{equation}
\lambda(\theta_j) \phi_j(a) u(a) (1 - R_j(a)) = e_j(a) (\delta R_j(a) + s)
\end{equation}

\end{document}